% Generated draft: requires mathematical and source review.
\documentclass[11pt]{article}
\usepackage[margin=25mm]{geometry}
\usepackage{fontspec}
\setmainfont{Latin Modern Roman}
\usepackage{amsmath,amssymb,mathtools,unicode-math}
\setmathfont{Latin Modern Math}
\usepackage{enumitem,needspace,xurl,hyperref}
\setlength{\parindent}{0pt}
\setlength{\parskip}{5pt}
\setlength{\emergencystretch}{3em}
\newcommand{\N}{\mathbb{N}}
\newcommand{\Z}{\mathbb{Z}}
\newcommand{\Q}{\mathbb{Q}}
\newcommand{\R}{\mathbb{R}}
\newcommand{\C}{\mathbb{C}}
\newcommand{\F}{\mathbb{F}}
\newcommand{\A}{\mathbb{A}}
\newcommand{\PP}{\mathbb{P}}
\newcommand{\E}{\mathbb{E}}
\newcommand{\eps}{\varepsilon}
\newcommand{\dd}{\,\mathrm{d}}
\begin{document}
\section*{UnsolvedMath ID: 30001315}
\textbf{Source problem code:} catalogue entry
\subsection{Definitions and mathematical statement}
The supplied record contains three conjectural assertions, although its normalized statement retains only the last two. The ambient space, connectedness, and boundary conventions are not explicit. The formulation below uses the customary interpretation of connected, orientable surfaces without boundary in Euclidean \(\R^3\); these conventions require confirmation against the source.

For such an embedded surface \(M\subset\R^3\), minimality means that its mean curvature vanishes identically, and completeness means completeness of its induced Riemannian metric. Proper embeddedness means that the inclusion \(M\hookrightarrow\R^3\) is proper. Write \(g(M)\) for its topological genus and \(e(M)\) for its number of topological ends, defined using the inverse system of components of complements of compact subsets. Finite total curvature means
\[
\int_M |K_M|\,\mathrm dA<\infty,
\]
where \(K_M\) is Gaussian curvature and \(\mathrm dA\) is induced area measure. Two surfaces are congruent if a Euclidean isometry of \(\R^3\) maps one onto the other.

\begin{enumerate}
\item \textbf{Genus-one classification, retained from the original statement.} Every properly embedded minimal surface of genus one is one of the ``known examples.'' The supplied record does not enumerate those examples or specify the equivalence relation for this classification. Consequently, this assertion cannot be made into a fully specified classification statement from the supplied evidence alone.

\item \textbf{One-ended genus-one uniqueness.} Let \(H_1\) denote the intended reference surface called the genus-one helicoid in the supplied record. The recorded question is whether, for every complete embedded minimal surface \(M\subset\R^3\),
\[
 g(M)=1\quad\text{and}\quad e(M)=1
 \quad\Longrightarrow\quad
 \exists\, T\in\operatorname{Isom}(\R^3)\ \text{such that}\ T(M)=H_1.
\]
The record does not specify a representative or scale normalization for \(H_1\). Since dilations preserve the stated hypotheses and congruence allows only isometries, the intended treatment of scale must be clarified. The displayed implication preserves the supplied word ``congruent'' and is an unresolved formulation pending that clarification.

\item \textbf{Finite-total-curvature end bound.} For every integer \(g\geq 0\) and every complete embedded minimal surface \(M\subset\R^3\) of finite total curvature and genus \(g\), is it true that
\[
 e(M)\leq g+2\,?
\]
The supplied literature assessment calls this the Hoffman--Meeks end bound. No properness hypothesis is added to this question.
\end{enumerate}

The background assessment also mentions uniqueness of a genus-\(g\) one-ended helicoid in general. This is broader than the explicit genus-one question, and no precise general-genus uniqueness formulation is supplied.
\subsection{Short English statement}
The record asks whether every complete embedded genus-one minimal surface with one end is congruent to the genus-one helicoid, and whether a complete embedded minimal surface of genus \(g\) and finite total curvature has at most \(g+2\) ends. Its original statement also proposes classification of properly embedded genus-one minimal surfaces by an unspecified collection of known examples. Ambient-space conventions, the classification list, and the scale convention for helicoid uniqueness require clarification. These are statement-only drafts requiring human mathematical review.
\subsection{Sources}
\begin{itemize}
\item Supplied local UnsolvedMath registry record, permanent ID \(30001315\), source problem code \texttt{catalogue entry}, title \emph{Genus-One Helicoid Classification and End Bounds}, in the Oberwolfach Reports Open Problems collection. No individual proposer is identified; the supplied proposed year is 2009.
\item Supplied report citation: \emph{Differentialgeometrie im Großen (2010)}, Oberwolfach Reports, pp.~1943--1972, DOI \texttt{10.4171/owr/2009/35}. Source URL and literature-source label ``Differential geometry report'': \url{https://doi.org/10.4171/owr/2009/35}. The supplied background describes Oberwolfach Reports as documenting mathematical workshops and their open-problem discussions. The citation's 2010 designation and the record's 2009 proposed year are retained without resolving their relationship.
\item Supplied literature source, labeled ``Genus-one existence result'': \url{https://doi.org/10.4007/annals.2009.169.347}. No authors or fuller bibliographic description are supplied.
\item Supplied literature source, labeled ``Minimal-surface conjecture survey'': \url{https://www.intlpress.com/site/pub/files/_fulltext/journals/sdg/2016/0021/0001/SDG-2016-0021-0001-a006.pdf}.
\item Literature assessment dated exactly \texttt{2026-08-21}: the supplied curation labels the record ``partially solved'' and reports that existence results cover important higher-genus examples, while uniqueness of a genus-\(g\) one-ended helicoid and the Hoffman--Meeks end bound remain open in general. This is supplied dated curation derived from the OWR material, not an independently verified literature assessment. The record expressly requests independent verification of time-sensitive status and cited claims.
\item The supplied statement status is \texttt{corrected\_verified}. Its accompanying verification note says the mathematical content was checked against the official Oberwolfach report and that notation, context, and question form were normalized without changing the source-level claim. This is a supplied verification claim; no independent source consultation or mathematical verification is asserted here.
\end{itemize}
\end{document}
